% SGA 3，第 XI 篇；根据英文重建源翻译。
% 范围：第 6 节全部，从标题、命题 6.1 至末注 6.12 及 Murre/Bourbaki 脚注；止于第 XII 篇之前。
% 权威来源：Polo--Gille Expo11.pdf，组件局部页 29--34；合卷页 751--756；印刷／原文页边页 170--179。
% Jacob Reinhold 的 CC BY 4.0 Markdown 谱系仅用于措辞比较；文字、公式、编号、注释及图均以 Polo--Gille PDF 为准。
% 精确续接／终点：第 XI 篇末尾，Expo11 局部页 34／合卷页 756／原文页边页 179；严格止于合卷页 757 的第 XII 篇之前。

\SGARefTarget{sga3:xi:section:6}\section*{6. 某些群概形同态的刚性性质\\及某些传送子的\\可表性\footnote{\SGARefLink{sga3:xi:section:6}{本节} 不使用第 \SGARefLink{sga3:xi:section:3}{3}、\SGARefLink{sga3:xi:section:4}{4}、\SGARefLink{sga3:xi:section:5}{5} 节的结果；它本应放在第 VI\textsubscript{B} 讲。}}\addcontentsline{toc}{section}{6. 某些群概形同态的刚性性质及某些传送子的可表性}

\medskip\noindent\SGARefTarget{sga3:xi:proposition:6:1}\textbf{命题 6.1。} --- \textit{设 \(S\) 为局部 Noether 预概形，\(H\) 为在 \(S\) 上有限型的交换 \(S\)-群预概形，\(E\) 为对乘法稳定的正整数集合。假设下列条件成立：}

\begin{enumerate}\item[(a)]\SGARefTarget{sga3:xi:proposition:6:1:item:a}\textit{对每个 \(n\in E\)，子群} \SGARefTarget{sga3:xi:display:6:1:n-torsion} \[
  {}_nH=\operatorname{Ker}\bigl(n\cdot\operatorname{id}_H\bigr)
\]\textit{在 \(S\) 上有限且平坦；}

\item[(b)]\SGARefTarget{sga3:xi:proposition:6:1:item:b}\textit{对每个 \(s\in S\)，各 \({}_nH_s\)（\(n\in E\)）构成的族在纤维 \(H_s\) 中概形稠密。\footnote{法文印刷文本此处为 \(H\)；逐纤维条件及其在证明中的使用要求为 \(H_s\)。}}\end{enumerate}

\textit{设 \(Y\) 为 \(H\) 的闭子预概形。则 \(S\) 的子函子 \(\prod_{H/S}(Y/H)\)（参照 \SGARefLink{sga3:viii:section:6}{第 VIII 篇，6}）可由 \(S\) 的闭子预概形 \(T\) 表示。若 \(S\) 为 Noether 预概形，则存在 \(n\in E\)，使得\footnote{印刷公式将右端的商写为 \((Y\cap{}_nH)/S\)。由于乘积函子在 \({}_nH/S\) 上取，且由证明的要求，此商应为 \((Y\cap{}_nH)/{}_nH\)。}} \SGARefTarget{sga3:xi:display:6:1:cross} \[
  \tag{x}
  \prod_{H/S}(Y/H)
  =
  \prod_{{}_nH/S}\bigl((Y\cap{}_nH)/{}_nH\bigr).
\]。

事实上，由 \SGARefLink{sga3:viii:theorem:6:4}{第 VIII 篇，6.4}，条件 \SGARefLink{sga3:xi:proposition:6:1:item:a}{(a)} 蕴含 \({}_nH\) 在 \(S\) 上有限且平坦，从而更是“本质自由”的。因此 \SGARefLink{sga3:xi:display:6:1:cross}{\((x)\)} 的右端可由 \(S\) 的闭子预概形 \(T_n\) 表示。\footnote{法文印刷文本此处为“\(T\) 的 \(T_n\)”；下一句将各 \(T_n\) 视为 \(S\) 的闭子预概形，而 \(T\) 仅在后面才定义为其稳定值。}按整除关系对 \(E\) 排序。各 \(T_n\) 构成 \(S\) 的闭子预概形的递减族；因此当 \(S\) 为 Noether 预概形时（现在可以这样假设），此族对充分大的 \(n\) 稳定。设 \(T\) 为各 \(T_n\) 的稳定值。断言 \(T\) 表示 \SGARefLink{sga3:xi:display:6:1:cross}{\((x)\)} 的左端，由此即完成证明。

只需证明下述结论。设 \(S'\) 为 \(S\) 上的预概形，且 \SGARefTarget{sga3:xi:display:6:1:base-change-condition} \[
  (Y\cap{}_nH)_{S'}=({}_nH)_{S'}
  \qquad\text{对每个 \(n\in E\)，}
\]，或等价地，对每个 \(n\in E\)，有 \(Y_{S'}\supset({}_nH)_{S'}\)。则 \(Y_{S'}=H_{S'}\)。这立即由 \SGARefLink{sga3:ix:lemma:4:4}{第 IX 篇，4.4} 得出：由条件 \SGARefLink{sga3:xi:proposition:6:1:item:a}{(a)} 和 \SGARefLink{sga3:xi:proposition:6:1:item:b}{(b)}，各 \(({}_nH)_{S'}\) 构成的族在 \(H_{S'}\) 中概形稠密。

\medskip\noindent\SGARefTarget{sga3:xi:theorem:6:2}\textbf{定理 6.2。} --- \textit{设 \(S\) 为预概形，\(H,E\) 如 \SGARefLink{sga3:xi:proposition:6:1}{6.1} 所述并满足条件 \SGARefLink{sga3:xi:proposition:6:1:item:a}{(a)} 和 \SGARefLink{sga3:xi:proposition:6:1:item:b}{(b)}，\(u:H\to G\) 为从 \(H\) 到在 \(S\) 上局部有限型的 \(S\)-群预概形的 \(S\)-群同态，\(K\) 为 \(G\) 的闭子群预概形。\footnote{法文印刷文本此处为 \(H\)；证明立即使用了 \(K\subset G\)，这正是传送子所要求的。}考虑 \(G\) 的子函子 \(\operatorname{Transp}_G(u,K)\)（参见 \SGARefLink{sga3:xi:proposition:2:5}{2.5}）。它可由 \(G\) 的闭子预概形表示。若 \(G\) 为 Noether 预概形，例如 \(S\) 为 Noether 预概形且 \(G\) 在 \(S\) 上有限型，则存在 \(n\in E\) 使得} \SGARefTarget{sga3:xi:display:6:2:transporter-stabilization} \[
  \operatorname{Transp}_G(u,K)
  =
  \operatorname{Transp}_G\bigl(u|_{{}_nH},K\bigr).
\]\textit{。最后，假设 \(G\) 在 \(S\) 上平滑，\(K\) 在 \(S\) 上平滑或为乘法型，且 \(H,E\) 满足下述比 \SGARefLink{sga3:xi:proposition:6:1:item:a}{(a)} 更强的条件：}

\begin{enumerate}\item[\(a'\)]\SGARefTarget{sga3:xi:theorem:6:2:item:a-prime}\textit{对每个 \(n\in E\)，\(H\) 的子群 \({}_nH\) 为乘法型。}\end{enumerate}

\textit{则 \(\operatorname{Transp}_G(u,K)\) 在 \(S\) 上平滑。}

考虑 \(G\)-群 \SGARefTarget{sga3:xi:display:6:2:hg} \[
  H_G=H\times_SG,
\]，它显然满足 \SGARefLink{sga3:xi:proposition:6:1}{6.1} 的各条件，其中将 \(S\) 换成 \(G\)，将 \(H\) 换成 \(H_G\)。设 \(Y\) 为 \(H_G\) 的闭子预概形，它是 \(K\subset G\) 在映射 \SGARefTarget{sga3:xi:display:6:2:conjugation-map} \[
  H_G=H\times_SG\longrightarrow G,\qquad
  (h,g)\longmapsto\operatorname{int}(g)\cdot u(h).
\] 下的逆像。则 \SGARefTarget{sga3:xi:display:6:2:product-identification} \[
  \operatorname{Transp}_G(u,K)
  =
  \prod_{H_G/G}(Y/H_G)
\]（参照 \SGARefLink{sga3:viii:examples:6:5:item:e:01}{第 VIII 篇例 6.5(e)}）。因此前几个断言由 \SGARefLink{sga3:xi:proposition:6:1}{6.1} 得出。此外，对每个拟紧开集 \(U\subset G\)，存在 \(n\in E\)，使两个函子 \SGARefTarget{sga3:xi:display:6:2:local-stabilization} \[
  \operatorname{Transp}_G(u,K)
  \qquad\text{且}\qquad
  \operatorname{Transp}_G\bigl(u|_{{}_nH},K\bigr)
\] 在 \(U\) 上的限制相同。因此，为验证 \SGARefLink{sga3:xi:theorem:6:2}{6.2} 的最后一个断言，可将 \(H\) 换成某个 \({}_nH\)，之后只需应用 \SGARefLink{sga3:xi:proposition:2:5}{2.5}，因为假设 \({}_nH\) 在 \(S\) 上为乘法型。

\medskip\noindent\SGARefTarget{sga3:xi:remark:6:3}\textbf{注 6.3。} --- 上述证明仅使用了 \SGARefLink{sga3:viii:theorem:6:4}{第 VIII 篇，6.4} 中非常初等的结果，最后一部分还使用了 \SGARefLink{sga3:xi:proposition:2:5}{2.5}；当 \(K\) 在 \(S\) 上平滑时，后者就是 \SGARefLink{sga3:ix:theorem:3:6}{第 IX 篇，3.6} 的无穷小结果，从而来自乘法型群的上同调消失。在最重要的情形中（参见 \SGARefLink{sga3:xi:unit:6:7}{6.7}），各整数 \(n\in E\) 甚至可以取为与 \(S\) 的剩余特征互素，即在 \(\mathscr O_S\) 中可逆。于是各 \({}_nH\) 在 \(S\) 上有限平展，所引用的上同调结果实际上是平凡的。因此在此情形中，\SGARefLink{sga3:xi:theorem:6:2}{6.2} 独立于乘法型群理论。

\medskip\noindent\SGARefTarget{sga3:xi:unit:6:4}\textbf{6.4。} --- 照例可知，\SGARefLink{sga3:xi:proposition:6:1}{6.1} 推广到下述情形。设 \(S'\) 为 \(S\)-预概形（不必局部 Noether），置 \(H'=H_{S'}\)，设 \(Y'\) 为 \(H'\) 的闭子预概形，并要求 \(Y'\to H'\) 有限呈示，即定义 \(Y'\) 的理想有限型。则 \(\prod_{H'/S'}(Y'/H')\) 可由 \(S'\) 的闭子预概形 \(T'\) 表示，态射 \(T'\to S'\) 有限呈示，且若 \(S'\) 拟紧，则存在 \(n\in E\)，使 \SGARefLink{sga3:xi:display:6:1:cross}{\((x)\)} 的类似结论成立。同样可知，只要浸入 \(K\to G\) 有限呈示，\SGARefLink{sga3:xi:theorem:6:2}{6.2} 的第一个断言仍成立，而无需假设 \(G\) 在 \(S\) 上局部有限型。

\medskip\noindent\SGARefTarget{sga3:xi:unit:6:5}\textbf{6.5。} --- 如 \SGARefLink{sga3:xi:remark:5:6}{5.6} 所预告，定理 \SGARefLink{sga3:xi:theorem:6:2}{6.2} 使得某些以另一方法在较强假设下对乘法型群证明的结果，可推广到满足上述 \SGARefLink{sga3:xi:theorem:6:2:item:a-prime}{\((a')\)} 和 \SGARefLink{sga3:xi:proposition:6:1:item:b}{(b)} 的群预概形。这适用于 \SGARefLink{sga3:xi:corollary:5:2}{5.2}、\SGARefLink{sga3:xi:corollary:5:3}{5.3} 的开头及 \SGARefLink{sga3:xi:corollary:5:4}{5.4} 和相应的“bis”变体，也适用于 \SGARefLink{sga3:xi:proposition:5:9}{5.9} 和 \SGARefLink{sga3:xi:corollary:5:10}{5.10}。还适用于第 \SGARefLink{sga3:ix:section:5}{第 IX 篇第 5 节}、6 节的结果，但 \SGARefLink{sga3:ix:theorem:6:8}{第 IX 篇，6.8} 除外。后一个断言甚至对剩余特征为 \(p>0\) 的离散赋值环的谱 \(S\) 上的阿贝尔概形同态 \(u:H\to G\) 已经不成立：\(\operatorname{Ker}u\) 的一般纤维可能是平凡群，而特殊纤维是非平凡的纯不可分的（radicial）群。

\medskip\noindent\SGARefTarget{sga3:xi:unit:6:6}\textbf{6.6。} --- 刚才给出了一个乘法型群具有而阿贝尔概形不具有的刚性性质的例子。另一个极其重要的例子是，当 \(H\) 为阿贝尔概形时，\SGARefLink{sga3:ix:theorem:3:6}{第 IX 篇，3.6} 中关于同态的无穷小延拓的存在定理不再成立。因此，域上的非零阿贝尔概形具有非平凡的无穷小变化，而乘法型群没有。这是以下事实的无穷小方面：阿贝尔簇有模理论，而且远未完备，而乘法型群的模理论则是空的。

这种无穷小差别还有如下“整体”方面。设 \(H\) 为局部 Noether 预概形 \(S\) 上的阿贝尔概形，\(G\) 为在 \(S\) 上局部有限型的交换 \(S\)-群预概形。可以证明 \(\operatorname{Hom}_{S\text{-}\mathrm{gr}}(H,G)\) 可由 \(S\) 上局部有限型的预概形表示。然而，与 \(H\) 为乘法型时的相应预概形不同，它在 \(S\) 上不平展，而仅非分歧。例如，若 \(S\) 为完备离散赋值环的谱，\(H,G\) 为 \(S\) 上的阿贝尔概形，则特殊纤维上可能存在同态 \(H_0\to G_0\)，它们并非来自一般纤维上的同态的“特殊化”。

\medskip\noindent\SGARefTarget{sga3:xi:unit:6:7}\textbf{6.7。} --- 当 \(H\) 为 \(S\) 上的阿贝尔概形，或更一般地，是这样的概形被环面扩张所得的群时，定理 \SGARefLink{sga3:xi:theorem:6:2}{6.2} 总适用。事实上，问题在 \(S\) 上是局部的，因此可假设存在与 \(S\) 的剩余特征互素的素数 \(\ell\)。取 \(E\) 为 \(\ell\) 的各次幂的集合，即足以满足条件 \SGARefLink{sga3:xi:theorem:6:2:item:a-prime}{\((a')\)} 和 \SGARefLink{sga3:xi:proposition:6:1:item:b}{(b)}。

与 \SGARefLink{sga3:xi:proposition:6:1}{6.1} 类似的论证给出 \SGARefLink{sga3:xi:theorem:6:8}{下述结果}。

\medskip\noindent\SGARefTarget{sga3:xi:theorem:6:8}\textbf{定理 6.8。} --- \textit{设 \(X\) 为 \(S\) 上平滑的预概形，具有非空连通几何纤维。对 \(X\) 的每个闭子预概形 \(Y\)，函子 \(\prod_{X/S}(Y/X)\) 可由 \(S\) 的闭子预概形 \(T\) 表示。若 \(Y\) 在 \(X\) 上有限呈示，则 \(T\) 在 \(S\) 上有限呈示。}

由于 \(f:X\to S\) 忠实平坦且局部有限呈示，它是 fpqc 拓扑的覆盖。另一方面，\(T=\prod_{X/S}(Y/X)\) 显然为 \(S\) 的 fpqc 子层。因此，\(T\) 可由 \(S\) 的闭子预概形表示，以及 \(T\) 是否在 \(S\) 上有限呈示这两个问题，在 \(S\) 上对 fpqc 拓扑都是局部的。作基变换 \(S'\to S\)，其中 \(S'=X\)，可归约到 \(X\) 在 \(S\) 上有截面 \(e\) 的情形。还可假设 \(S\) 仿射，特别地拟紧。于是有：

\medskip\noindent\SGARefTarget{sga3:xi:corollary:6:9}\textbf{推论 6.9。} --- \textit{在 \SGARefLink{sga3:xi:theorem:6:8}{6.8} 的假设下，假设 \(S\) 拟紧，\(X\) 在 \(S\) 上有截面 \(e\)。对每个整数 \(n\geq0\)，以 \(X_n\) 表示 \(X\) 中截面 \(e\) 的 \(n\) 阶无穷小邻域。假设 \(Y\) 在 \(X\) 上有限呈示。则存在整数 \(n\geq0\)，使得} \SGARefTarget{sga3:xi:display:6:9:stabilization} \[
  \prod_{X/S}(Y/X)
  =
  \prod_{X_n/S}(Y_n/X_n),
  \qquad Y_n=Y\cap X_n.
\]。

当 \(Y\) 在 \(X\) 上有限呈示时，此推论通过 \SGARefLink{sga3:viii:theorem:6:4}{第 VIII 篇，6.4} 蕴含 \SGARefLink{sga3:xi:theorem:6:8}{6.8}。事实上，由于 \(X\) 在 \(S\) 上平滑，\(X_n\) 在 \(S\) 上有限且局部自由，从而更在 \(S\) 上是 \SGARefLink{sga3:viii:definition:6:1}{第 VIII 篇，6.1} 意义下“本质自由”的。因此 \(\prod_{X_n/S}(Y_n/X_n)\) 可由 \(S\) 的闭子预概形表示。此外，所引用结果的证明，或归约到 Noether 情形，均立即表明此闭子预概形在 \(S\) 上有限呈示。

先在 \(S\) 为 Noether 预概形时证明 \SGARefLink{sga3:xi:corollary:6:9}{6.9}，从而证明 \SGARefLink{sga3:xi:theorem:6:8}{6.8}。置 \SGARefTarget{sga3:xi:display:6:9:tn} \[
  T_n=\prod_{X_n/S}(Y_n/X_n).
\]。各 \(T_n\) 构成 \(S\) 的闭子预概形的递减序列；由于 \(S\) 为 Noether 预概形，此序列稳定。设 \SGARefTarget{sga3:xi:display:6:9:r} \[
  R=\bigcap_{n\geq0}\prod_{X_n/S}(Y_n/X_n)
\] 为充分大的 \(n\) 对应的共同值。显然 \(T\subset R\)，还须说明 \(R\subset T\)。作基变换 \(R\to S\) 后，归约到 \(R=S\) 的情形，即对每个 \(n\)，有 \(Y_n=X_n\)；等价地，对每个 \(n\)，有 \(Y\supset X_n\)。这时须证明 \(T=S\)，即 \(Y=X\)。

对所有 \(n\) 成立的包含 \(Y\supset X_n\)，以及 \(X\) 局部 Noether 这一事实，蕴含在 \(e(S)\) 的每个点的某个邻域内，\(Y\) 为 \(X\) 的带诱导结构的开子预概形。因此，存在 \(X\) 的带诱导结构的开子预概形 \(U\)，使得 \(e(S)\subset U\subset Y\)。由 \SGARefLink{sga3:ix:corollary:4:3}{第 IX 篇，4.3}，由于 \(X/S\) 的纤维整，\(U\) 在 \(X\) 中概形稠密。\(Y\) 为包含 \(U\) 的闭子预概形，因此 \(Y=X\)。这就证明了 Noether 情形中的 \SGARefLink{sga3:xi:corollary:6:9}{6.9}，从而证明了 \SGARefLink{sga3:xi:theorem:6:8}{6.8}。

一般情形归约到前述情形。对每个 \(s\in S\)，存在 \(s\) 的仿射开邻域 \(U\) 及 \(e(s)\) 的仿射开邻域 \(V\)，使得 \(f(V)\subset U\)。则 \(f(V)\) 为包含于 \(U\) 的 \(s\) 的开邻域。选取包含于 \(e^{-1}(V)\cap f(V)\) 的 \(s\) 的仿射开邻域 \(S'\)，置 \(X'=V\cap f^{-1}(S')\)。则 \(X'\)、\(S'\) 分别为 \(X\)、\(S\) 的仿射开子预概形，且 \(X'/S'\) 有截面。由 \SGARefLink{sga3:xi:theorem:6:8}{6.8} 和 \SGARefLink{sga3:xi:corollary:6:9}{6.9} 的局部性，可以假设 \(S=S'\)。

断言 \SGARefTarget{sga3:xi:display:6:9:open-reduction} \[
  \prod_{X/S}(Y/X)=\prod_{X'/S'}(Y'/X'),
  \qquad Y'=Y\cap X'.
\]。事实上，由 \SGARefLink{sga3:ix:corollary:4:6}{第 IX 篇，4.6}，至少当 \(X\) 在 \(S\) 上拟分离，从而 \(X'\) 在 \(X\) 中逆紧时，\(X'\) 在 \(X\) 中概形稠密。实际上，容易证明 \SGARefLink{sga3:ix:corollary:4:6}{第 IX 篇，4.6} 在没有逆紧性假设时仍成立。同样，在每个基变换 \(S_1\to S\) 后，\(X'\times_SS_1\) 在 \(X\times_SS_1\) 中概形稠密，所断言的等式立即随之成立。因此归约到 \(X=X'\) 的情形，从而可以假设 \(S\) 和 \(X\) 均仿射。

现在设 \(X=\operatorname{Spec}(B)\)，\(J\) 为 \(B\) 中定义 \(Y\) 的理想。理想 \(J\) 是其有限型子理想的滤过余极限。因此，\(Y\) 是 \(X\) 中在 \(S\) 上有限呈示的闭子预概形 \(Y_i\) 的交，且 \SGARefTarget{sga3:xi:display:6:9:intersection} \[
  \prod_{X/S}(Y/X)
  =
  \bigcap_i\prod_{X/S}(Y_i/X).
\]。所以为证明 \SGARefLink{sga3:xi:theorem:6:8}{6.8}，只需处理 \(Y\) 在 \(X\) 上有限呈示的情形。还须在 \(S\) 和 \(X\) 仿射时证明 \SGARefLink{sga3:xi:corollary:6:9}{6.9}。但此时 \(S\) 上的 \(X\) 和 \(Y\) 可经基变换 \(S\to S_0\) 从 Noether 预概形 \(S_0\) 上的类似情形 \(X_0,Y_0\) 得到。这样就归约到已处理的 \(S\) 为 Noether 预概形的情形，完成了 \SGARefLink{sga3:xi:theorem:6:8}{6.8} 和 \SGARefLink{sga3:xi:corollary:6:9}{6.9} 的证明。

\medskip\noindent\SGARefTarget{sga3:xi:corollary:6:10}\textbf{推论 6.10。} --- \textit{设 \(X\) 为平滑、有限呈示且具有连通纤维的 \(S\)-群预概形，\(Y\) 为 \(S\) 上有限呈示的群预概形，\(i:Y\to X\) 为 \(S\)-群预概形的单态射，从而 \(Y\) 为 \(X\) 的子群。则 \(\prod_{X/S}(Y/X)\) 可由 \(S\) 的有限呈示闭子预概形表示。若 \(S\) 拟紧，对每个 \(n\geq0\)，以 \(X_n\) 表示 \(X\) 的单位截面的 \(n\) 阶无穷小邻域，并置 \(Y_n=X_n\cap Y\)。则对所有充分大的 \(n\)，} \SGARefTarget{sga3:xi:display:6:10:stabilization} \[
  \prod_{X/S}(Y/X)=\prod_{X_n/S}(Y_n/X_n).
\]。

证明本质上就是 \SGARefLink{sga3:xi:corollary:6:9}{6.9} 的证明。\(X\) 和 \(Y\) 的单位截面诱导双射浸入 \(S\to X_n\) 和 \(S\to Y_n\)，因而诱导同构 \SGARefTarget{sga3:xi:display:6:10:reductions} \[
  S_{\mathrm{red}}\xrightarrow{\sim}(X_n)_{\mathrm{red}},
  \qquad
  S_{\mathrm{red}}\xrightarrow{\sim}(Y_n)_{\mathrm{red}}.
\]。由此可知，单射 \(Y_n\to X_n\) 为固有态射。由于它是有限呈示单态射，因此为闭浸入。前面已使用过的 \SGARefLink{sga3:viii:theorem:6:4}{第 VIII 篇，6.4} 现表明 \(\prod_{X_n/S}(Y_n/X_n)\) 可由 \(S\) 的闭子预概形表示，且在 \(S\) 上有限呈示。还须在 \(S\) 仿射的附加假设下证明 \SGARefLink{sga3:xi:corollary:6:10}{6.10} 的最后一个断言。

再次归约到 \(S\) 为 Noether 预概形的情形。用 \SGARefLink{sga3:xi:corollary:6:9}{6.9} 的证明中的记号，只需证明 \(R=T\)，或等价地，对每个 \(n\) 有 \(Y\supset X_n\) 蕴含 \(Y=X\)。此假设蕴含 \(i:Y\to X\) 在 \(S\) 上 \(Y\) 的单位截面的各点处平展，因此 \(Y\) 在这些点处于 \(S\) 上平滑。故由 \(Y\) 在 \(S\) 上平滑的点构成的开子预概形 \(Y'\)，是 \(Y\) 的带诱导结构的开子群。由 \SGARefLink{sga3:x:lemma:3:5}{第 X 篇，3.5}，态射 \(Y'\to X\) 为平展单态射，因而为开浸入。由于 \(X\) 的纤维连通，且代数群的每个开子群也闭，此开浸入为满射，因而为同构。所以 \(Y'=X\)，从而更有 \(Y=X\)，这就完成了 \SGARefLink{sga3:xi:corollary:6:10}{6.10} 的证明。

与 \SGARefLink{sga3:viii:examples:6:5}{第 VIII 篇，6.5} 中一样论证，得到：

\medskip\noindent\SGARefTarget{sga3:xi:corollary:6:11}\textbf{推论 6.11。} --- \textit{设 \(G\) 为在 \(S\) 上局部有限型且拟分离的 \(S\)-群预概形，\(H\) 为 \(S\) 上平滑、有限呈示且具有连通纤维的群预概形，\(i:H\to G\) 为 \(S\)-群的单态射，从而 \(H\) 为 \(G\) 的子群。则：}

\begin{enumerate}\item[(a)]\SGARefTarget{sga3:xi:corollary:6:11:item:a}\textit{\(\operatorname{Centr}_G(H)\) 和 \(\operatorname{Norm}_G(H)\) 可由 \(G\) 中在 \(G\) 上有限呈示的闭子预概形表示。类似地，对 \(S\)-群预概形之间的每个有限呈示单态射 \(j:K\to G\)，\(\operatorname{Transp}_G(H,K)\) 可由 \(G\) 中在 \(G\) 上有限呈示的闭子预概形表示。}

\item[(b)]\SGARefTarget{sga3:xi:corollary:6:11:item:b}\textit{若 \(G\) 拟紧，则在 \SGARefLink{sga3:xi:corollary:6:11:item:a}{(a)} 的各情形中，都存在整数 \(n\geq0\)，使得若以 \(H_n\) 表示 \(H\) 的单位截面的 \(n\) 阶无穷小邻域（在传送子情形中类似地定义 \(K_n\)），则} \SGARefTarget{sga3:xi:display:6:11:stabilization} \[
  \begin{aligned}
  \operatorname{Centr}_G(H)
    &=\operatorname{Centr}_G(H_n),\\
  \operatorname{Norm}_G(H)
    &=\operatorname{Norm}_G(H_n),\\
  \operatorname{Transp}_G(H,K)
    &=\operatorname{Transp}_G(H_n,K)
     =\operatorname{Transp}_G(H_n,K_n).
  \end{aligned}
\]。\end{enumerate}

在基预概形 \(G\) 上，对群预概形 \SGARefTarget{sga3:xi:display:6:11:hg} \[
  X=H_G=H\times_SG.
\] 应用 \SGARefLink{sga3:xi:corollary:6:10}{6.10}。对于中心化子，取 \(Y\) 为 \((G\times_SG)_G\) 的对角子群在从 \(X\) 到 \((G\times_SG)_G\) 的适当 \(G\)-群同态下的逆像。对于传送子，取 \(K_G\) 在从 \(X\) 到 \(G_G\) 的适当 \(G\)-群同态下的逆像。正规化子情形通过取 \(K=H\) 归约到传送子情形。关于 \(G\) 的假设保证 \(H\to G\) 有限呈示，从而 \(Y\to X\) 有限呈示。对于中心化子，同一假设保证 \(G\times_SG\) 的对角子群在 \(G\times_SG\) 上有限呈示；再次得出 \(Y\) 在 \(X\) 上有限呈示。

\medskip\noindent\SGARefTarget{sga3:xi:remark:6:12}\textbf{注 6.12。} --- 利用 EGA VI 中较为细致的结果，\footnote{另参见 J.\,P. Murre，“非分歧函子的表示．应用〔Representation of unramified functors. Applications〕”，\emph{Séminaire Bourbaki}，第\ 294 号（1965 年 5 月），定理 3，第\ 13 页。}可以证明：若 \(X\) 为 \(S\) 上平坦且有限呈示的预概形，并且在 \(S\) 上固有或具有非空连通几何纤维，则对 \(X\) 的每个在 \(S\) 上有限呈示的闭子预概形 \(Y\)，函子 \(\prod_{X/S}(Y/X)\) 可由 \(S\) 中在 \(S\) 上有限呈示的闭子预概形表示。类似地，若 \(X\) 为平坦、有限呈示且具有连通纤维的 \(S\)-群预概形，\(i:Y\to X\) 为 \(S\)-群的单态射，其中 \(Y\) 为有限呈示 \(S\)-群预概形，则 \(\prod_{X/S}(Y/X)\) 可由 \(S\) 中在 \(S\) 上有限呈示的闭子预概形表示。特别地，将假设“\(H\) 在 \(S\) 上平滑”换为“\(H\) 在 \(S\) 上平坦”后，\SGARefLink{sga3:xi:corollary:6:11:item:a}{6.11(a)} 仍成立。


